\section{期望}

\subsection{数学期望}
\begin{definition}
	设$f$是概率空间$(X,\mathscr{F},P)$上的随机变量。若$f$在$X$上的积分存在，则称$f$的\gls{MathematicalExpectation}（简称为期望）存在，并将：
	\begin{equation*}
		\operatorname{E}(f)\coloneq\int_{X}f(x)\dif P
	\end{equation*}
	称为$f$的期望。若$f$在$X$上可积，则称$f$的期望是有限的。根据\cref{theo:ProductMeasurable}，随机向量的期望即为各分量期望按原顺序排列成的向量，随机矩阵的期望即为各元素期望按原顺序排成的矩阵。
\end{definition}
\begin{property}\label{prop:Expectation}
	设$f$是概率空间$(X,\mathscr{F},P)$上在$X$上积分存在的随机变量。期望有如下性质：
	\begin{enumerate}
		\item 对任何$(\mathbb{R},\mathcal{B}(\mathbb{R}))$上的可测函数$g$，$g\circ f$是$(X,\mathscr{F},P)$上的可测函数，只要：
		\begin{equation*}
			\operatorname{E}(g\circ f),\quad\int_{\mathbb{R}}g\dif Pf^{-1}
		\end{equation*}
		之一有意义，则二者一定相等。特别的，当$f$为连续型随机变量时，上式化作：
		\begin{equation*}
			\operatorname{E}(g\circ f),\quad\int_{\mathbb{R}}g(x)\frac{\dif Pf^{-1}}{\dif\lambda}(x)\dif\lambda
		\end{equation*}
		其中$\lambda$为L测度；当$f$为离散型随机变量时，令$D=\{x_n\}$，含义与随机变量分类时的含义相同，则上式化作：
		\begin{equation*}
			\operatorname{E}(g\circ f),\quad\sum_{n=1}^{+\infty}g(x_n)P(f=x_n)
		\end{equation*}
		\item 若$\seq{f}{n}$是$(X,\mathscr{F},P)$上相互独立的随机变量，$\seq{g}{n}$是$(\mathbb{R},\mathcal{B}(\mathbb{R}))$上的Borel函数，只要：
		\begin{equation*}
			\operatorname{E}\left[\prod_{i=1}^{n}g_i(f_i)\right],\quad\prod_{i=1}^{n}\operatorname{E}[g_i(f_i)]
		\end{equation*}
		之一有意义，二者一定相等。
	\end{enumerate}
\end{property}
\begin{proof}
	(1)由\cref{prop:MeasurableMapping}(2)可知$g\circ f$是$(X,\mathscr{F},P)$到$(\mathbb{R},\mathcal{B}(\mathbb{R}))$上的可测函数。注意到：
	\begin{equation*}
		\operatorname{E}(g\circ f)=\int_{X}g[f(x)]\dif P=\int_{f^{-1}(\mathbb{R})}g[f(x)]\dif P
	\end{equation*}
	由\cref{theo:IntBySubstitution}即可得出一般结论。根据\cref{lem:IntChangeOfMeasure}可得出连续型随机变量时的结论。对于离散型随机变量，由\cref{prop:MeasurableFunction}(2)、\cref{prop:SigmaField}(4)、\cref{theo:MeasurableCountableIntegral}、\cref{prop:MeasurableIntegral}(1)和积分的定义（$g(x)$限制在$x_n$上是简单函数，若$g(x_n)<0$还需用到\cref{prop:MeasurableIntegral}(6)）可得：
	\begin{align*}
		&\int_{\mathbb{R}}g(x)\dif Pf^{-1}=\int_{(\mathbb{R}\setminus D)\bigcup\left(\underset{n=1}{\overset{+\infty}{\cup}}\{x_n\}\right)}g(x)\dif Pf^{-1} \\
		=&\int_{\mathbb{R}\setminus D}g(x)\dif Pf^{-1}+\sum_{n=1}^{+\infty}\int_{\{x_n\}}g(x)\dif Pf^{-1}=\sum_{n=1}^{+\infty}g(x_n)P(f=x_n)
	\end{align*}\par
	(2)由\cref{prop:MeasurableFunction}(11)（可以证明$g(\seq{x}{n})=x_1x_2\cdots x_n$在$\mathbb{R}^{n}$上连续）和\cref{prop:Independent}(5.d)即可得到。
\end{proof}
\begin{note}
	很多教材在上述定理中写的并不是$g$对pushforward measure$Pf^{-1}$进行积分，而是对$f$的分布函数$F$进行积分。$F$是一个测度吗？$F$与$Pf^{-1}$等价吗？显然并不是。由L-S测度处的讨论可以看出$F$可唯一确定$Pf^{-1}$，所以可使用$F$来代表$Pf^{-1}$。
\end{note}

\subsection{条件期望}
通常的期望$\operatorname{E}(f)$描述随机变量$f$在整个概率空间上的均值，但在实际问题中，我们还需要研究在某些条件成立时$f$的均值。例如，对于事件$A\in\mathscr{F}$，可以研究$f$在$A$发生条件下的平均值。然而，如果可能给出的条件有很多，我们并不希望对每一个条件分别定义一个彼此独立的数值，而希望构造一个新的随机变量，使得条件发生变化时，这个随机变量能够对应$f$在这一条件下的均值。当我们希望把给定条件后的均值组织成一个随机变量时，需要同时处理一整套的条件，这一整套条件自然应该构成一个$\sigma$域。
\begin{lemma}\label{lem:ConditionalExpectation}
	设$f$为概率空间$(X,\mathscr{F},P)$上在$X$上积分存在的随机变量。由\cref{prop:MeasurableIntegral}(3)，对任意的$A\in\mathscr{F}$，令：
	\begin{equation*}
		\varphi(A)=\int_{A}f(x)\dif P
	\end{equation*}
	则：
	\begin{enumerate}
		\item $\varphi$是$\mathscr{F}$上的符号测度；
		\item $\varphi\ll P$。
	\end{enumerate}
\end{lemma}
\begin{proof}
	(1)因为$P(\varnothing)=0$，由\cref{prop:MeasurableIntegral}(1)可得$\varphi(\varnothing)=0$。根据\cref{theo:MeasurableCountableIntegral}可知对互不相交的$\{A_n\}\subseteq\mathscr{F}$有：
	\begin{equation*}
		\varphi\left(\underset{n=1}{\overset{+\infty}{\cup}}A_n\right)=\int_{\underset{n=1}{\overset{+\infty}{\cup}}A_n}f(x)\dif P=\sum_{n=1}^{+\infty}\left[\int_{A_n}f(x)\dif P\right]=\sum_{n=1}^{+\infty}\varphi(A_n)
	\end{equation*}
	所以$\varphi$是$\mathscr{F}$上的符号测度。\par
	(2)由\cref{prop:MeasurableIntegral}(1)直接可得。
\end{proof}
\begin{definition}
	设$f$为概率空间$(X,\mathscr{F},P)$上在$X$上积分存在的随机变量，$\mathscr{A}$是$\mathscr{F}$的一个子$\sigma$域。对任意的$A\in\mathscr{F}$，令：
	\begin{equation*}
		\varphi(A)=\int_{A}f(x)\dif P
	\end{equation*}
	若存在$(X,\mathscr{A})$上的可测函数$g$满足对任意的$A\in\mathscr{A}$有：
	\begin{equation*}
		\int_{A}g(x)\dif P=\int_{A}f(x)\dif P
	\end{equation*}
	则称$g$为$f$关于$\mathscr{A}$的\gls{ConditionalExpectation}，记作$\operatorname{E}(f\mid\mathscr{A})$。当$\mathscr{A}=\sigma(h)$且$h$是$(X,\mathscr{F})$到可测空间 $(Y,\mathscr{B})$的可测映射时，称$\operatorname{E}(f\mid h)\coloneq\operatorname{E}[f|\sigma(h)]$为$f$关于$h$的条件期望（由\cref{lem:PreimageSigmaField}可知$\sigma(h)$为$\mathscr{F}$的子$\sigma$域）。当给定多个子$\sigma$域$\mathscr{A}_1,\mathscr{A}_2,\dots,\mathscr{A}_n\subseteq\mathscr{F}$时，约定：
	\begin{equation*}
		\operatorname{E}(f\mid\mathscr{A}_1,\mathscr{A}_2,\dots,\mathscr{A}_n)\coloneq\operatorname{E}\left[f\mid\sigma\left(\underset{i=1}{\overset{n}{\cup}}\mathscr{A}_i\right)\right]
	\end{equation*}
\end{definition}
\begin{note}
	因为概率是$\sigma$有限测度，由\cref{prop:Measure}(4)、\cref{lem:ConditionalExpectation}(1)、\cref{prop:SignedMeasure}(7)、\cref{lem:ConditionalExpectation}(2)和\cref{theo:RadonNikodym}可知存在$(X,\mathscr{A},P)$上a.e.于$(X,\mathscr{A},P)$的意义下唯一的在$X$上积分存在的可测函数$g$满足：
	\begin{equation*}
		\forall\;A\in\mathscr{A},\;\varphi(A)=\int_{A}g(x)\dif P
	\end{equation*}
	所以条件期望必然存在且在a.e.于$(X,\mathscr{A},P)$的意义下唯一。
\end{note}
作为一个随机变量，$\operatorname{E}(f\mid\mathscr{A})$在某一点$x\in X$处的取值究竟表示什么？因为$\operatorname{E}(f\mid\mathscr{A})$是几乎处处意义下的，所以我们也没必要关心它在零测集上的取值。进一步地，概率空间中任意一族两两不交且具有正概率的可测集合至多可数。若$\{A_i:i\in I\}\subseteq\mathscr{A}$两两不交且$P(A_i)>0$，对任意$n\in\mathbb{N}^{+}$令：
\begin{equation*}
	I_n=\left\{i\in I:P(A_i)\geqslant\frac{1}{n}\right\}
\end{equation*}
因为这些集合两两不交且$P(X)=1$，由\cref{prop:Measure}(1)(3)（单调性）可得$|I_n|\leqslant n$，从而每个$I_n$均为有限集。而：
\begin{equation*}
	I=\underset{n=1}{\overset{+\infty}{\cup}}I_n
\end{equation*}
由\cref{theo:CountableUnionCountable}可知$I$至多可数。
\begin{theorem}\label{theo:PartitionMeasurableFunction}
	设$\{A_i:i\in I\}$是$X$的至多可数划分，即：
	\begin{equation*}
		X=\underset{i\in I}{\cup}A_i,\quad A_i\cap A_j=\varnothing\;(i\neq j)
	\end{equation*}
	$\mathscr{A}=\sigma(\{A_i:i\in I\})$，则$f$是$(X,\mathscr{A})$上的可测函数当且仅当存在$\{a_i:i\in I\}\subseteq\mathbb{R}$使得：
	\begin{equation*}
		f=\sum_{i\in I}a_iI(x\in A_i)
	\end{equation*}
\end{theorem}
\begin{proof}
	\textbf{(1)必要性：}任取$i\in I$，若存在$x_1,x_2\in A_i$使得$f(x_1)=\alpha<\beta=f(x_2)$，取$c=\dfrac{\alpha+\beta}{2}$。因为$f$关于$\mathscr{A}$可测，所以$f^{-1}((-\infty,c))\in\mathscr{A}$。由$\{A_i:i\in I\}$的定义可知：
	\begin{equation*}
		\mathscr{A}=\left\{\underset{i\in J}{\cup}A_i:J\subseteq I\right\}
	\end{equation*}
	所以$\mathscr{A}$中的集合要么包含$A_i$要么与$A_i$不相交。因为：
	\begin{equation*}
		x_1\in f^{-1}((-\infty,c)),\quad x_2\notin f^{-1}((-\infty,c))
	\end{equation*}
	所以此时$f$不可测，矛盾，于是$f$在每个$A_i$上均为常数。\par
	\textbf{(2)充分性：}由可测函数的定义即可证得。
\end{proof}
\begin{theorem}\label{theo:ConditionalExpectationPartition}
	设$f$是概率空间$(X,\mathscr{F},P)$上积分存在的随机变量，$\{A_i:i\in I\}\subseteq\mathscr{F}$是$X$的至多可数划分，$\mathscr{A}=\sigma(\{A_i:i\in I\})$，
	则：
	\begin{equation*}
		\operatorname{E}(f\mid\mathscr{A})(x)=\sum_{\substack{i\in I\\P(A_i)>0}}\frac{\int_{A_i}f\dif P}{P(A_i)}I(x\in A_i)\;\text{a.e.于$(X,\mathscr{A},P)$}
	\end{equation*}
	其中，在满足$P(A_i)=0$的集合$A_i$上可以任意定义$\operatorname{E}(f\mid\mathscr{A})$的取值。
\end{theorem}
\begin{proof}
	记：
	\begin{equation*}
		g(x)=\sum_{\substack{i\in I\\P(A_i)>0}}\frac{\displaystyle\int_{A_i}f\dif P}{P(A_i)}I(x\in A_i)
	\end{equation*}
	并在满足$P(A_i)=0$的集合$A_i$上任意定义$g$的取值，由\cref{theo:PartitionMeasurableFunction}可知$g$关于$\mathscr{A}$可测。\par
	任取$A\in\mathscr{A}$。由：
	\begin{equation*}
		\mathscr{A}=\left\{\underset{i\in J}{\cup}A_i:J\subseteq I\right\}
	\end{equation*}
	可知存在$J\subseteq I$使得：
	\begin{equation*}
		A=\underset{i\in J}{\cup}A_i
	\end{equation*}
	根据\cref{theo:MeasurableCountableIntegral}和\cref{prop:MeasurableIntegral}(1)可得：
	\begin{align*}
		\int_{A}g(x)\dif P=\sum_{\substack{i\in J\\P(A_i)>0}}\int_{A_i}g(x)\dif P=\sum_{i\in J}\int_{A_i}f(x)\dif P=\int_{A}f(x)\dif P
	\end{align*}
	由条件期望的定义可知结论成立。
\end{proof}
\cref{theo:ConditionalExpectationPartition}给出了一个较为一般且直观的情形，在这一情形下可以明确说明条件期望作为随机变量时每一点取值的含义。当$\mathscr{A}$由至多可数划分$\{A_i:i\in I\}$生成时，对于$P(A_i)>0$的$x\in A_i$，条件期望在$x$处的取值恰好为$f$在事件$A_i$发生条件下的均值：
\begin{equation*}
	\operatorname{E}(f\mid\mathscr{A})(x)=\frac{1}{P(A_i)}\int_{A_i}f\dif P
\end{equation*}
而对于满足$P(A_i)=0$的划分块，由$\{A_i:i\in I\}$至多可数以及\cref{prop:Measure}(3)（次可列可加性）可知它们的并集仍为零测集。因此，除去一个零测集以后，条件期望在每一点的取值都可以具体理解为该点所属条件下$f$的均值，前面所希望构造的“随着条件变化而给出相应条件均值的随机变量”在这里得到了具体实现。
\begin{property}\label{prop:ConditionalExpectation}
	设$f,g$是概率空间$(X,\mathscr{F},P)$上在$X$上积分存在的随机变量，$\mathscr{A},\mathscr{B}$是$\mathscr{F}$的子$\sigma$域，则：
	\begin{enumerate}
		\item $f$关于$\mathscr{A}$可测的充要条件为$\operatorname{E}(f\mid\mathscr{A})=f\;$a.e.于$(X,\mathscr{A},P)$；
		\item 若$f$与$\mathscr{A}$独立，则$\operatorname{E}(f\mid\mathscr{A})=\operatorname{E}(f)\;$a.e.于$(X,\mathscr{A},P)$；
		\item 若$\mathscr{A}\subseteq\mathscr{B}$，则$\operatorname{E}[\operatorname{E}(f\mid\mathscr{A})|\mathscr{B}]=\operatorname{E}(f\mid\mathscr{A})\;$a.e.于$(X,\mathscr{B},P)$，$\operatorname{E}[\operatorname{E}(f\mid\mathscr{B})|\mathscr{A}]=\operatorname{E}(f\mid\mathscr{A})\;$a.e.于$(X,\mathscr{A},P)$，$\operatorname{E}[\operatorname{E}(f\mid\mathscr{A})]=\operatorname{E}(f)$；
		\item 若$f\leqslant g\;$a.e.于$(X,\mathscr{A},P)$，则$\operatorname{E}(f\mid\mathscr{A})\leqslant\operatorname{E}(g|\mathscr{A})\;$a.e.于$(X,\mathscr{A},P)$；
		\item 对任意的$\alpha,\beta\in\mathbb{R}$，若$\alpha\operatorname{E}(f)+\beta\operatorname{E}(g)$有意义，则：
		\begin{equation*}
			\operatorname{E}(\alpha f+\beta g|\mathscr{A})=\alpha\operatorname{E}(f\mid\mathscr{A})+	\beta\operatorname{E}(g|\mathscr{A})
		\end{equation*}
		a.e.于$(X,\mathscr{A},P)$；
		\item 若$0\leqslant f_n\uparrow f\;$a.e.于$(X,\mathscr{A},P)$，则$0\leqslant\operatorname{E}(f_n|\mathscr{A})\uparrow \operatorname{E}(f\mid\mathscr{A})\;$a.e.于$(X,\mathscr{A},P)$；
		\item 若$f_n\geqslant0\;$a.e.于$(X,\mathscr{F},P)$，则$\operatorname{E}\left(\varliminf\limits_{n\to+\infty}f_n\Big|\mathscr{A}\right)\leqslant\varliminf\limits_{n\to+\infty}\operatorname{E}(f_n|\mathscr{A})\;$a.e.于$(X,\mathscr{A},P)$；
		\item 若$|f_n|\leqslant g\in L_1(X,\mathscr{A},P)$且$f_n\overset{\text{a.e.}}{\longrightarrow}f$，则$\lim\limits_{n\to+\infty}\operatorname{E}(f_n|\mathscr{A})=\operatorname{E}(f\mid\mathscr{A})\;$a.e.于$(X,\mathscr{A},P)$。
	\end{enumerate}
\end{property}
\begin{proof}
	(1)必要性由条件期望的定义即可得到，充分性由条件期望的定义和\cref{prop:MeasurableFunction}(9)即可得到。\par
	(2)因为$f$与$\mathscr{A}$独立，由\cref{prop:Independent}(3)可知对任意的$A\in\mathscr{A}$有$f$与$I_A$独立。对任意的$A\in\mathscr{A}$，根据\cref{prop:MeasurableIntegral}(3)、\cref{prop:SigmaField}(4)、\cref{theo:MeasurableCountableIntegral}、\cref{prop:Expectation}(2)和\cref{prop:MeasurableIntegral}(6)可得：
	\begin{align*}
		&\int_{A}f(x)\dif P=\int_{A}f(x)\dif P+\int_{X\setminus A}0\dif P=\operatorname{E}(fI_A) \\
		=&\operatorname{E}(f)\operatorname{E}(I_A)=\operatorname{E}(f)P(A)=\operatorname{E}(f)\int_{A}1\dif P=\int_{A}\operatorname{E}(f)\dif P
	\end{align*}
	于是结论成立。\par
	(3)由$\operatorname{E}(f\mid\mathscr{A})$的定义可知$\operatorname{E}(f\mid\mathscr{A})$是$\mathscr{A}$可测的，因为$\mathscr{A}\subseteq\mathscr{B}$，所以$\operatorname{E}(f\mid\mathscr{A})$是$\mathscr{B}$可测的。由(1)可得$\operatorname{E}[\operatorname{E}(f\mid\mathscr{A})|\mathscr{B}]=\operatorname{E}(f\mid\mathscr{A})\;$a.e.于$(X,\mathscr{B},P)$。\par
	由$\operatorname{E}[\operatorname{E}(f\mid\mathscr{B})|\mathscr{A}]$的定义可知它是$\mathscr{A}$可测的，于是对任意的$A\in\mathscr{A}$，根据条件期望的定义可得：
	\begin{equation*}
		\int_{A}\operatorname{E}[\operatorname{E}(f\mid\mathscr{B})|\mathscr{A}]\dif P=\int_{A}\operatorname{E}(f\mid\mathscr{B})\dif P=\int_{A}f(x)\dif P
	\end{equation*}
	即$\operatorname{E}(f\mid\mathscr{A})=\operatorname{E}[\operatorname{E}(f\mid\mathscr{B})|\mathscr{A}]\;$a.e.于$(X,\mathscr{A},P)$。\par
	因为$X\in\mathscr{A}$，由条件期望的定义可得：
	\begin{equation*}
		\operatorname{E}[\operatorname{E}(f\mid\mathscr{A})]=\int_{X}\operatorname{E}(f\mid\mathscr{A})\dif P=\int_{X}f(x)\dif P=\operatorname{E}(f)
	\end{equation*}\par
	(4)对任意的$A\in\mathscr{A}$，由\cref{prop:MeasurableIntegral}(7)可得：
	\begin{equation*}
		\int_{A}\operatorname{E}(f\mid\mathscr{A})\dif P=\int_{A}f(x)\dif P\leqslant\int_{A}g(x)\dif P=\int_{A}\operatorname{E}(g|\mathscr{A})\dif P
	\end{equation*}
	由条件期望的定义可知$\operatorname{E}(f\mid\mathscr{A}),\operatorname{E}(g|\mathscr{A})$关于$\mathscr{A}$可测，根据\cref{prop:MeasurableFunction}(3)可得$E=\{\operatorname{E}(f\mid\mathscr{A})>\operatorname{E}(g|\mathscr{A})\}\in\mathscr{A}$，所以：
	\begin{equation*}
		\int_{E}\operatorname{E}(f\mid\mathscr{A})\dif P\leqslant\int_{E}\operatorname{E}(g|\mathscr{A})\dif P
	\end{equation*}
	由\cref{prop:MeasurableIntegral}(7)可得：
	\begin{equation*}
		\int_{E}\operatorname{E}(f\mid\mathscr{A})\dif P\geqslant\int_{E}\operatorname{E}(g|\mathscr{A})\dif P
	\end{equation*}
	所以：
	\begin{equation*}
		\int_{E}\operatorname{E}(f\mid\mathscr{A})\dif P=\int_{E}\operatorname{E}(g|\mathscr{A})\dif P
	\end{equation*}
	根据\cref{prop:MeasurableIntegral}(6)可知：
	\begin{equation*}
		\int_{E}[\operatorname{E}(f\mid\mathscr{A})-\operatorname{E}(g|\mathscr{A})]\dif P=0
	\end{equation*}
	由\cref{prop:NonnegativeMeasurableIntegral}(9)可得$P(E)=0$，即$\operatorname{E}(f\mid\mathscr{A})\leqslant\operatorname{E}(g|\mathscr{A})\;$a.e.于$(X,\mathscr{A},P)$。\par
	(5)因为$\alpha\operatorname{E}(f)+\beta\operatorname{E}(g)$存在，由\cref{prop:MeasurableIntegral}(6)可知$\operatorname{E}(\alpha f+\beta g)$存在，于是$\operatorname{E}(\alpha f+\beta g|\mathscr{A})$存在。对任意的$A\in\mathscr{A}$，有：
	\begin{equation*}
		\int_{A}[\alpha f(x)+\beta g(x)]\dif P=\alpha\int_{A}f(x)\dif P+\beta\int_{A}g(x)\dif P
	\end{equation*}
	否则由\cref{prop:MeasurableIntegral}(6)可知右式不存在，即$\alpha\int_{A}f(x)\dif P$和$\beta\int_{A}g(x)\dif P$二者异号无穷大，分类讨论可知此时与$\alpha\operatorname{E}(f)+\beta\operatorname{E}(g)$存在矛盾（对$f^+,f^-,g^+,g^-$在$X$上和$A$上的积分进行分析）。\par
	根据\cref{prop:MeasurableIntegral}(6)，对任意的$A\in\mathscr{A}$有：
	\begin{align*}
		\int_{A}\operatorname{E}(\alpha f+\beta g|\mathscr{A})\dif P&=\int_{A}[\alpha f(x)+\beta g(x)]\dif P=\alpha\int_{A}f(x)\dif P+\beta\int_{A}g(x)\dif P \\
		&=\alpha\int_{A}\operatorname{E}(f\mid\mathscr{A})\dif P+\beta\int_{A}\operatorname{E}(g|\mathscr{A})\dif P \\
		&=\int_{A}\alpha\operatorname{E}(f\mid\mathscr{A})\dif P+\int_{A}\beta\operatorname{E}(g|\mathscr{A})\dif P \\
		&=\int_{A}[\alpha\operatorname{E}(f\mid\mathscr{A})+\beta\operatorname{E}(g|\mathscr{A})]\dif P
	\end{align*}
	由条件期望的定义可得$\operatorname{E}(\alpha f+\beta g|\mathscr{A})=\alpha\operatorname{E}(f\mid\mathscr{A})+\beta\operatorname{E}(g|\mathscr{A})\;$a.e.于$(X,\mathscr{A},P)$。\par
	(6)由(4)可得$0\leqslant\operatorname{E}(f_n|\mathscr{A})\;$a.e.于$(X,\mathscr{A},P)$且$\operatorname{E}(f_n|\mathscr{A})\uparrow\;$a.e.于$(X,\mathscr{A},P)$，根据\cref{prop:RSeq}(7)可知$\lim\limits_{n\to+\infty}\operatorname{E}(f_n|\mathscr{A})$存在，由\cref{prop:MeasurableFunction}(6)可知$\lim\limits_{n\to+\infty}\operatorname{E}(f_n|\mathscr{A})$关于$\mathscr{A}$可测。根据\cref{theo:LeviTheorem}可知对任意的$A\in\mathscr{A}$有：
	\begin{equation*}
		\int_{A}\left[\lim_{n\to+\infty}\operatorname{E}(f_n|\mathscr{A})\right]\dif P=\lim_{n\to+\infty}\left[\int_{A}\operatorname{E}(f_n|\mathscr{A})\dif P\right]=\lim_{n\to+\infty}\left[\int_{A}f_n(x)\dif P\right]=\int_{A}f(x)\dif P
	\end{equation*}
	由条件期望的定义可得：
	\begin{equation*}
		\lim_{n\to+\infty}\operatorname{E}(f_n|\mathscr{A})=\operatorname{E}(f\mid\mathscr{A})
	\end{equation*}
	a.e.于$(X,\mathscr{A},P)$。根据\cref{theo:CountableUnionCountable}和\cref{prop:Measure}(3)（次可列可加性）即可得出结论。\par
	(7)因为$f_n\geqslant0\;$a.e.于$(X,\mathscr{A},P)$，所以：
	\begin{equation*}
		0\leqslant\inf_{i\geqslant n}f_i\uparrow\lim_{n\to+\infty}\inf\limits_{i\geqslant n}f_i\;\text{a.e.于}(X,\mathscr{A},P)
	\end{equation*}
	由(4)和\cref{prop:Measure}(3)（次可列可加性）可得对任意的$n\in\mathbb{N}^+$有：
	\begin{equation*}
		\operatorname{E}\left(\inf_{i\geqslant n}f_i\Big|\mathscr{A}\right)\leqslant\inf_{i\geqslant n}\operatorname{E}(f_i|\mathscr{A})
	\end{equation*}
	a.e.于$(X,\mathscr{A},P)$。根据(6)、\cref{prop:RSeq}(6)、\cref{prop:SigmaField}(3)和\cref{prop:Measure}(3)（次有限可加性）可得：
	\begin{align*}
		&\operatorname{E}\left(\varliminf_{n\to+\infty}f_n\Big|\mathscr{A}\right)=\operatorname{E}\left(\lim_{n\to+\infty}\inf_{i\geqslant n}f_i\Big|\mathscr{A}\right)\leqslant\lim_{n\to+\infty}\operatorname{E}\left(\inf_{i\geqslant n}f_i\Big|\mathscr{A}\right) \\
		\leqslant&\lim_{n\to+\infty}\inf_{i\geqslant n}\operatorname{E}(f_i|\mathscr{A})=\varliminf_{n\to+\infty}\operatorname{E}(f_n|\mathscr{A})
	\end{align*}
	a.e.于$(X,\mathscr{A},P)$。\par
	(8)由(7)类似证明\cref{cor:FatouLemma}，然后类似\cref{theo:DominatedConvergenceTheorem}的证明即可得出结论。
\end{proof}
\begin{theorem}\label{theo:ConditionalExpectationPullOut}
	设$f,g$是概率空间$(X,\mathscr{F},P)$上的随机变量，$f$和$fg$在$X$上积分存在，$g$关于$\mathscr{F}$的子$\sigma$域$\mathscr{A}$可测，则：
	\begin{equation*}
		\operatorname{E}(fg|\mathscr{A})=g\operatorname{E}(f\mid\mathscr{A})
	\end{equation*}
	a.e.于$(X,\mathscr{A},P)$。
\end{theorem}
\begin{proof}
	当$f\geqslant0\;$a.e.于$(X,\mathscr{A},P)$上时，使用典型方法。\par
	对任意的$A,B\in\mathscr{A}$，当$g(x)=I(x\in B)$时，由\cref{prop:SigmaField}(2)(4)、\cref{prop:MeasurableIntegral}(3)和\cref{theo:MeasurableCountableIntegral}可得：
	\begin{align*}
		&\int_{A}\operatorname{E}(fg|\mathscr{A})\dif P=\int_{A}f(x)g(x)\dif P=\int_{A}f(x)I(x\in B)\dif P \\
		=&\int_{A\cap B}f(x)\dif P+\int_{A\setminus B}0\dif P=\int_{A\cap B}f(x)\dif P=\int_{A\cap B}\operatorname{E}(f\mid\mathscr{A})\dif P \\
		=&\int_{A\cap B}\operatorname{E}(f\mid\mathscr{A})I(x\in B)\dif P+\int_{A\setminus B}\operatorname{E}(f\mid\mathscr{A})I(x\in B)\dif P=\int_{A}\operatorname{E}(f\mid\mathscr{A})I(x\in B)\dif P
	\end{align*}
	所以结论对指示函数成立。\par
	当$g$是非负简单函数时，设$X$可表示为有限个互不相交的集合$\seq{E}{n}\in\mathscr{A}$的并，$\seq{a}{n}\geqslant0$，$g=\sum\limits_{i=1}^{n}a_iI(x\in E_i)$。对任意的$i=1,2,\dots,n$，由\cref{prop:MeasurableIntegral}(8)可知$\sum\limits_{i=1}^{n}\operatorname{E}[fa_iI(x\in E_i)]$存在，根据\cref{prop:ConditionalExpectation}(5)、\cref{prop:MeasurableFunction}(5.b)、\cref{prop:MeasurableIntegral}(8)、指示函数时的情况和\cref{prop:NonnegativeMeasurableIntegral}(10)可得：
	\begin{align*}
		&\operatorname{E}(fg|\mathscr{A})=\operatorname{E}\left[\sum_{i=1}^{n}fa_iI(x\in E_i)\Big|\mathscr{A}\right]=\sum_{i=1}^{n}\operatorname{E}[fa_iI(x\in E_i)|\mathscr{A}] \\
		=&\sum_{i=1}^{n}I(x\in E_i)\operatorname{E}(fa_i|\mathscr{A})=\sum_{i=1}^{n}a_iI(x\in E_i)\operatorname{E}(f\mid\mathscr{A})=g\operatorname{E}(f\mid\mathscr{A})
	\end{align*}
	所以结论对非负简单函数成立。\par
	当$g$是非负随机变量时，由\cref{prop:MeasurableFunction}(8)可知存在非负简单函数列$\{g_n\}$满足$g_n\uparrow g$，于是$fg_n\uparrow fg$且$fg_n\geqslant0\;$a.e.于$(X,\mathscr{A},P)$。根据\cref{prop:RSeq}(8.c)、\cref{prop:ConditionalExpectation}(6)和非负简单函数时的情况可得：
	\begin{equation*}
		\operatorname{E}(fg|\mathscr{A})=\operatorname{E}\left(\lim_{n\to+\infty}fg_n\Big|\mathscr{A}\right)=\lim_{n\to+\infty}\operatorname{E}(fg_n|\mathscr{A})=\lim_{n\to+\infty}g_n\operatorname{E}(f\mid\mathscr{A})=g\operatorname{E}(f\mid\mathscr{A})
	\end{equation*}
	所以结论对非负随机变量成立。\par
	当$g$是一般随机变量时，$g=g^+-g^-$。由\cref{prop:MeasurableIntegral}(8)可得：
	\begin{equation*}
		\operatorname{E}(fg^+)-\operatorname{E}(fg^-)=\operatorname{E}(fg)
	\end{equation*}
	存在，根据\cref{prop:ConditionalExpectation}(5)和非负随机变量时的情况可得：
	\begin{align*}
		&\operatorname{E}(fg|\mathscr{A})=\operatorname{E}(fg^+-fg^-|\mathscr{A})=\operatorname{E}(fg^+|\mathscr{A})-\operatorname{E}(fg^-|\mathscr{A}) \\
		=&g^+\operatorname{E}(f\mid\mathscr{A})-g^-\operatorname{E}(f\mid\mathscr{A})=g\operatorname{E}(f\mid\mathscr{A})
	\end{align*}
	所以结论对一般随机变量成立。\par
	当$f$为一般随机变量时，$(fg)^+=f^+g^++f^-g^-,\;(fg)^-=f^+g^-+f^-g^+$。因为$\operatorname{E}(fg)$存在，所以$\operatorname{E}[(fg)^+]$和$\operatorname{E}[(fg)^-]$不同时为正无穷。假设$\operatorname{E}[(fg)^+]$不为正无穷，$\operatorname{E}[(f g)^-]$不为正无穷的情况可类似讨论。因为$f^+g^+,f^-g^-\leqslant(fg)^+$，由\cref{prop:NonnegativeMeasurableIntegral}(6)可得：
	\begin{equation*}
		\operatorname{E}(f^+g^+),\operatorname{E}(f^-g^-)\leqslant\operatorname{E}[(fg)^+]<+\infty
	\end{equation*}
	根据\cref{prop:MeasurableIntegral}(6)可得：
	\begin{align*}
		&\operatorname{E}(f^+g)-\operatorname{E}(f^-g)=\operatorname{E}(f^+g^+-f^+g^-)-\operatorname{E}(f^-g^+-f^-g^-) \\
		=&\operatorname{E}(f^+g^+)-\operatorname{E}(f^+g^-)-\operatorname{E}(f^-g^+)+\operatorname{E}(f^-g^-)=\operatorname{E}(f^+g^++f^-g^-)-\operatorname{E}(f^+g^-+f^-g^+)
	\end{align*}
	存在。由\cref{prop:ConditionalExpectation}(5)和$f\geqslant0\;$a.e.于$(X,\mathscr{A},P)$时的情况可得：
	\begin{align*}
		&\operatorname{E}(fg|\mathscr{A})=\operatorname{E}(f^+g-f^-g|\mathscr{A})=\operatorname{E}(f^+g|\mathscr{A})-\operatorname{E}(f^-g|\mathscr{A}) \\
		=&g[\operatorname{E}(f^+|\mathscr{A})-\operatorname{E}(f^-|\mathscr{A})]=g\operatorname{E}(f\mid\mathscr{A})\qedhere
	\end{align*}
\end{proof}

\subsection{条件概率}
\begin{definition}
	设$f$为概率空间$(X,\mathscr{F},P)$上的可测映射，$\mathscr{A}$是$\mathscr{F}$的一个子$\sigma$域。对任意的$A\in\mathscr{F}$，由\cref{prop:SimpleFunction}(1)可知$I_A$在$X$上积分存在，于是分别称：
	\begin{equation*}
		P(A|\mathscr{A})\coloneq\operatorname{E}(I_A|\mathscr{A}),\quad P(A|f)\coloneq P(A|\sigma(f))
	\end{equation*}
	为事件$A$关于$\mathscr{A}$的\gls{ConditionalProbability}和事件$A$关于$f$的条件概率。
\end{definition}
这一构造与前面对条件期望的解释完全一致。当$\mathscr{A}$由至多可数划分$\{A_i:i\in I\}$生成时，由\cref{theo:ConditionalExpectationPartition}、\cref{prop:SigmaField}(2)(4)和\cref{theo:MeasurableCountableIntegral}可知对于$P(A_i)>0$的$x\in A_i$有：
\begin{equation*}
	P(A|\mathscr{A})(x)=\frac{1}{P(A_i)}\int_{A_i}I_A\dif P=\frac{P(A\cap A_i)}{P(A_i)}
\end{equation*}
因此，关于$\sigma$域的条件概率作为随机变量在$x$处的取值正是当已知$x$所属的条件$A_i$发生时事件$A$发生的概率。换言之，关于$\sigma$域的条件概率并不是一个单独的数值，而是将随条件变化的一整套条件概率统一组织成了一个随机变量。
\begin{definition}
	设$(X,\mathscr{F},P)$为概率空间，$A,B\in\mathscr{F}$且$P(B)>0$。因为$P(A\mid\sigma(B))$为$\sigma(B)$可测函数，由\cref{theo:PartitionMeasurableFunction}可知存在$a_1,a_2\in\mathbb{R}$使得\footnote{这里的a.e.来源于条件期望的a.e.。}：
	\begin{equation*}
		P(A\mid\sigma(B))=a_1I_B+a_2I_{B^c}\;\text{a.e.于$(X,\sigma(B),P)$}
	\end{equation*}
	称$a_1$为事件$A$在事件$B$发生条件下的条件概率，记为$P(A\mid B)\coloneq a_1$。
\end{definition}
\begin{property}\label{prop:ConditionalProbability}
	设$(X,\mathscr{F},P)$为概率空间，则条件概率具有如下性质：
	\begin{enumerate}
		\item(条件概率公式) 对任意$A,B\in\mathscr{F}$且$P(B)>0$，有：
		\begin{equation*}
			P(A\mid B)=\frac{P(A\cap B)}{P(B)}
		\end{equation*}
		\item(全概率公式) 对任意$A\in\mathscr{F}$以及$\mathscr{F}$的任意子$\sigma$域$\mathscr{A}$，有：
		\begin{equation*}
			P(A)=\operatorname{E}[P(A\mid\mathscr{A})]
		\end{equation*}
		\item(贝叶斯公式) 对任意$A,B\in\mathscr{F}$且$P(A)>0,\;P(B)>0$，有：
		\begin{equation*}
			P(B\mid A)=\frac{P(A\mid B)P(B)}{\operatorname{E}[P(A\mid\sigma(B))]}
		\end{equation*}
	\end{enumerate}
\end{property}
\begin{proof}
	(1)由$P(A\mid B)$的定义可知：
	\begin{equation*}
		P(A\mid\sigma(B))=P(A\mid B)I_B+a_2I_{B^c}\;\text{a.e.于$(X,\sigma(B),P)$}
	\end{equation*}
	其中$a_2\in\mathbb{R}$。由\cref{prop:MeasurableIntegral}(8)(6)、\cref{prop:SigmaField}(2)(4)和\cref{theo:MeasurableCountableIntegral}可得：
	\begin{gather*}
		\int_{B}P(A\mid\sigma(B))\dif P=\int_{B}[P(A\mid B)I_B(x)+a_2I_{B^c}(x)]\dif P=P(A\mid B)P(B) \\
		\int_{B}I_A(x)\dif P=\int_{A\cap B}I_A(x)\dif P+\int_{B\backslash A}I_A(x)\dif P=P(A\cap B)
	\end{gather*}
	根据$P(A\mid\sigma(B))$的定义可得：
	\begin{equation*}
		P(A\mid B)P(B)=P(A\cap B)
	\end{equation*}
	由于$P(B)>0$，因此：
	\begin{equation*}
		P(A\mid B)=\frac{P(A\cap B)}{P(B)}
	\end{equation*}\par
	(2)对于$\mathscr{F}$任意的子$\sigma$域$\mathscr{A}$，因为$X\in\mathscr{A}$，根据条件期望的定义和\cref{theo:MeasurableCountableIntegral}可知：
	\begin{equation*}
		\int_{X}\operatorname{E}(I_A\mid\mathscr{A})\dif P=\int_{X}I_A(x)\dif P=P(A)
	\end{equation*}
	又因为：
	\begin{equation*}
		P(A\mid\mathscr{A})=\operatorname{E}(I_A\mid\mathscr{A})
	\end{equation*}
	所以：
	\begin{equation*}
		\operatorname{E}[P(A\mid\mathscr{A})]=P(A)
	\end{equation*}\par
	(3)由(1)可得：
	\begin{equation*}
		P(B\mid A)P(A)=P(A\cap B)=P(A\mid B)P(B)
	\end{equation*}
	因为$P(A)>0$，根据(2)可知：
	\begin{equation*}
		P(B\mid A)=\frac{P(A\mid B)P(B)}{P(A)}=\frac{P(A\mid B)P(B)}{\operatorname{E}[P(A\mid\sigma(B))]}\qedhere
	\end{equation*}
\end{proof}
\subsubsection{条件概率分布}
为了描述给定随机变量$g$的取值后随机变量$f$的概率分布，自然需要一个同时依赖于$g$的取值与$f$的可测集合的映射：固定$g$的取值时，该映射关于$f$的可测集合应构成一个概率测度；而固定$f$的可测集合时，对应的概率应当是$g$取值的可测函数，从而使得该概率仍构成一个可测随机变量。
\begin{definition}
	设$(X,\mathscr{F})$和$(Y,\mathscr{A})$是可测空间。若映射$K:X\times\mathscr{A}\to[0,1]$满足：
	\begin{enumerate}
		\item 对任意的$x\in X$，$K(x,\cdot)$是$(Y,\mathscr{A})$上的概率测度；
		\item 对任意的$A\in\mathscr{A}$，$K(\cdot,A)$是$(X,\mathscr{F})$上的可测函数。
	\end{enumerate}
	则称$K$为从$(X,\mathscr{F})$到$(Y,\mathscr{A})$的\gls{ProbabilityKernel}。
\end{definition}
\begin{definition}
	设$f,g$分别是从概率空间$(X,\mathscr{F},P)$到可测空间$(Y,\mathscr{A})$和$(Z,\mathscr{B})$的可测映射。若从$(Z,\mathscr{B})$到$(Y,\mathscr{A})$的概率核$K_{f\mid g}$满足对任意的$A\in\mathscr{A}$和$B\in\mathscr{B}$有：
	\begin{equation*}
		P(\{f\in A\}\cap\{g\in B\})=\int_{B}K_{f\mid g}(z,A)\dif Pg^{-1}
	\end{equation*}
	则称$K_{f\mid g}$为$f$关于$g$的\gls{ConditionalProbabilityKernel}。
\end{definition}
对固定的$z\in Z$，$K_{f\mid g}(z,\cdot)$是$(Y,\mathscr{A})$上的概率测度，因此可以将其理解为当$g$的取值为$z$时$f$的一个候选分布。条件概率核的定义进一步要求，对任意$A\in\mathscr{A}$和$B\in\mathscr{B}$，事件$\{f\in A\}\cap\{g\in B\}$的概率能够通过在$g\in B$的各个可能取值上，对$K_{f\mid g}(z,A)$按照$g$的分布$Pg^{-1}$进行积分得到。这个积分关系说明$K_{f\mid g}(z,A)$具有给定$g=z$时$f\in A$的概率的含义，而下面的性质进一步表明它与条件概率$P(f\in A\mid g)$实际上是一致的。
\begin{property}\label{prop:ConditionalProbabilityKernel}
	设$f,g$分别是从概率空间$(X,\mathscr{F},P)$到可测空间$(Y,\mathscr{A})$和$(Z,\mathscr{B})$的可测映射。概率核$K_{f\mid g}$是$f$关于$g$的条件概率核当且仅当对任意的$A\in\mathscr{A}$有：
	\begin{equation*}
		K_{f\mid g}[g(\omega),A]=P(f\in A\mid g)(\omega)\;\text{a.e.于$(X,\sigma(g),P)$}
	\end{equation*}
\end{property}
\begin{proof}
	设$K_{f\mid g}$是$f$关于$g$的条件概率核。任取$A\in\mathscr{A}$，由概率核的定义和\cref{prop:MeasurableMapping}(2)可知$K_{f\mid g}[g(\omega),A]$关于$\sigma(g)$可测。\par
	\textbf{(1)必要性：}任取$B\in\mathscr{B}$，由\cref{theo:IntBySubstitution}、\cref{prop:SigmaField}(2)和\cref{theo:MeasurableCountableIntegral}可得：
	\begin{align*}
		&\int_{g^{-1}(B)}K_{f\mid g}[g(\omega),A]\dif P=\int_{B}K_{f\mid g}(z,A)\dif Pg^{-1} \\
		=&P(\{f\in A\}\cap\{g\in B\})=\int_{g^{-1}(B)}I(f\in A)\dif P
	\end{align*}
	由条件概率的定义即可得出结论。\par
	\textbf{(2)充分性：}任取$A\in\mathscr{A}$和$B\in\mathscr{B}$，由\cref{theo:IntBySubstitution}、\cref{prop:MeasurableIntegral}(8)、\cref{prop:SigmaField}(2)和\cref{theo:MeasurableCountableIntegral}可得：
	\begin{align*}
		&\int_{B}K_{f\mid g}(z,A)\dif Pg^{-1}=\int_{g^{-1}(B)}K_{f\mid g}[g(\omega),A]\dif P \\
		=&\int_{g^{-1}(B)}P(f\in A\mid g)(\omega)\dif P=\int_{g^{-1}(B)}\operatorname{E}(\{f\in A\}\mid g)\dif P \\
		=&\int_{g^{-1}(B)}I(f\in A)\dif P=P(\{f\in A\}\cap\{g\in B\})
	\end{align*}
	充分性得证。
\end{proof}
\subsubsection{条件概率密度}
\begin{definition}
	设$f,g$分别是从概率空间$(X,\mathscr{F},P)$到可测空间$(Y,\mathscr{A})$和$(Z,\mathscr{B})$的可测映射，$K_{f\mid g}$是$f$关于$g$的条件概率核，$\mu$是$(Y,\mathscr{A})$上的$\sigma$有限测度。若$(Y\times Z,\mathscr{A}\times\mathscr{B})$上的非负可测函数$p_{f\mid g}$满足：存在$N\in\mathscr{B}$使得$Pg^{-1}(N)=0$，并且对任意的$z\in Z\setminus N$和$A\in\mathscr{A}$有：
	\begin{equation*}
		K_{f\mid g}(z,A)=\int_{A}p_{f\mid g}(y\mid z)\dif\mu
	\end{equation*}
	则称$p_{f\mid g}(y\mid z)$为$f$关于$g$相对于$\mu$的\gls{ConditionalProbabilityDensityFunction}。
\end{definition}
这里允许排除一个满足$Pg^{-1}(N)=0$的集合$N$，是因为条件概率核本身只在$g$的分布意义下几乎处处确定（\cref{prop:MeasurableIntegral}(8)），因此其对应的条件概率密度也只需要对$Pg^{-1}$几乎处处的$z$定义一致。
\begin{property}\label{prop:ConditionalProbabilityDensityBasic}
	设$f,g$分别是从概率空间$(X,\mathscr{F},P)$到可测空间$(Y,\mathscr{A})$和$(Z,\mathscr{B})$的可测映射，$\mu$是$(Y,\mathscr{A})$上的$\sigma$有限测度，则：
	\begin{enumerate}
		\item 若$p_{f\mid g}$为$f$关于$g$相对于$\mu$的条件概率密度函数，则对任意的$A\in\mathscr{A}$有：
		\begin{equation*}
			P(f\in A\mid g)(\omega)=K_{f\mid g}[g(\omega),A]=\int_{A}p_{f\mid g}[y\mid g(\omega)]\dif\mu\;\text{a.e.于$(X,\sigma(g),P)$}
		\end{equation*}
		并且对任意的$B\in\mathscr{B}$有：
		\begin{equation*}
			P(\{f\in A\}\cap\{g\in B\})=\int_{B}\left[\int_{A}p_{f\mid g}(y\mid z)\dif\mu\right]\dif Pg^{-1}
		\end{equation*}
		\item $f$关于$g$相对于$\mu$的条件概率密度函数存在当且仅当$P(f,g)^{-1}\ll\mu\times Pg^{-1}$。当该绝对连续条件成立时，$f$关于$g$的条件概率核也存在，并且：
		\begin{gather*}
			\forall\;z\in Z,\;A\in\mathscr{A},\;K_{f\mid g}(z,A)=\int_Ap_{f\mid g}(y\mid z)\dif\mu \\
			p_{f\mid g}(y\mid z)
			=\dfrac{\dif P(f,g)^{-1}}{\dif(\mu\times Pg^{-1})}(y,z)\;\text{a.e.于$(Y\times Z,\mathscr{A}\times\mathscr{B},\mu\times Pg^{-1})$}
		\end{gather*}
	\end{enumerate}
\end{property}
\begin{proof}
	(1)因为$Pg^{-1}(N)=0$，所以$P(g\in N)=0$。由\cref{prop:ConditionalProbabilityKernel}、\cref{prop:SigmaField}(3)和\cref{prop:Measure}(3)（次有限可加性）可得：
	\begin{equation*}
		P(f\in A\mid g)(\omega)=K_{f\mid g}[g(\omega),A]=\int_{A}p_{f\mid g}[y\mid g(\omega)]\dif\mu\;\text{a.e.于$(X,\sigma(g),P)$}
	\end{equation*}
	根据\cref{prop:MeasurableIntegral}(8)可得：
	\begin{equation*}
		P(\{f\in A\}\cap\{g\in B\})=\int_{B}K_{f\mid g}(z,A)\dif Pg^{-1}=\int_{B}\left[\int_{A}p_{f\mid g}(y\mid z)\dif\mu\right]\dif Pg^{-1}
	\end{equation*}\par
	(2)由\cref{theo:Fubini}(1)可知$\mu\times Pg^{-1}$是$(Y\times Z,\mathscr{A}\times\mathscr{B})$上的$\sigma$有限测度。\par
	\textbf{必要性：}设$p_{f\mid g}$为$f$关于$g$相对于$\mu$的条件概率密度函数。对任意的$C\in\mathscr{A}\times\mathscr{B}$，定义：
	\begin{equation*}
		\varphi(C)\coloneq\int_{C}p_{f\mid g}(y\mid z)\dif(\mu\times Pg^{-1})
	\end{equation*}
	由\cref{theo:Fubini}(2)和\cref{prop:NonnegativeMeasurableIntegral}(2)可知对任意的$A\in\mathscr{A}$和$B\in\mathscr{B}$有：
	\begin{align*}
		\varphi(A\times B)&=\int_{A\times B}p_{f\mid g}(y\mid z)\dif(\mu\times Pg^{-1})=\int_{B}\left[\int_{A}p_{f\mid g}(y\mid z)\dif\mu\right]\dif Pg^{-1} \\
		&=\int_{B}K_{f\mid g}(z,A)\dif Pg^{-1}\geqslant0
	\end{align*}
	同时：
	\begin{equation*}
		\varphi(Y\times Z)=\int_{Z}K_{f\mid g}(z,Y)\dif Pg^{-1}=\int_{Z}1\dif Pg^{-1}=Pg^{-1}(Z)=P(X)=1
	\end{equation*}
	再结合\cref{theo:MeasurableCountableIntegral}可知$\varphi$是$(Y\times Z,\mathscr{A}\times\mathscr{B})$上的概率测度。由(1)可知对任意的$A\in\mathscr{A},\;B\in\mathscr{B}$有：
	\begin{equation*}
		\varphi(A\times B)=P(\{f\in A\}\cap\{g\in B\})=P(f,g)^{-1}(A\times B)
	\end{equation*}
	根据\cref{prop:FiniteDimensionalCylinderSet}(1)、\cref{theo:SetNecessarilySet1}和\cref{lem:PiExtensionMeasure}可得$\varphi=P(f,g)^{-1}$。由\cref{prop:MeasurableIntegral}(1)和Radon-Nikodym导数的定义可知：
	\begin{equation*}
		P(f,g)^{-1}\ll\mu\times Pg^{-1},\quad
		p_{f\mid g}=\dfrac{\dif P(f,g)^{-1}}{\dif(\mu\times Pg^{-1})}\;\text{a.e.于$(Y\times Z,\mathscr{A}\times\mathscr{B},\mu\times Pg^{-1})$}
	\end{equation*}\par
	\textbf{充分性：}设$P(f,g)^{-1}\ll\mu\times Pg^{-1}$，由\cref{theo:RadonNikodym}，取非负可测函数：
	\begin{equation*}
		q(y,z)=\dfrac{\dif P(f,g)^{-1}}{\dif(\mu\times Pg^{-1})}(y,z),\quad r(z)=\int_{Y}q(y,z)\dif\mu
	\end{equation*}
	由\cref{theo:Fubini}(2)可知$r$关于$\mathscr{B}$可测。任取$B\in\mathscr{B}$，由\cref{theo:Fubini}(2)可得：
	\begin{align*}
		&\int_{B}r(z)\dif Pg^{-1}=\int_{Y\times B}q(y,z)\dif(\mu\times Pg^{-1}) \\
		=&P(f,g)^{-1}(Y\times B)=Pg^{-1}(B)=\int_{B}1\dif Pg^{-1}
	\end{align*}
	由\cref{theo:RadonNikodym}的唯一性可知$r(z)=1\;$a.e.于$(Z,\mathscr{B},Pg^{-1})$（$Pg^{-1}$关于$Pg^{-1}$的Radon-Nikodym导数）。记：
	\begin{equation*}
		N=\{z\in Z:r(z)\ne1\}
	\end{equation*}
	则$Pg^{-1}(N)=0$，根据\cref{prop:MeasurableFunction}(3)可知$N\in\mathscr{B}$。因为$P(f,g)^{-1}$是概率测度且$P(f,g)^{-1}\ll\mu\times Pg^{-1}$，所以$\mu$不是零测度。又因为$\mu$是$\sigma$有限测度，所以存在$E\in\mathscr{A}$使得$0<\mu(E)<+\infty$。定义：
	\begin{equation*}
		\rho(y)\coloneq\frac{I(y\in E)}{\mu(E)},\quad
		p_{f\mid g}(y\mid z)\coloneq
		\begin{cases}
			q(y,z),&z\in Z\setminus N \\
			\rho(y),&z\in N
		\end{cases}
	\end{equation*}
	由\cref{prop:MeasurableFunction}(1.a)可验证得到$p_{f\mid g}$关于$\mathscr{A}\times\mathscr{B}$可测，根据\cref{prop:SigmaField}(4)、\cref{theo:MeasurableCountableIntegral}和\cref{prop:MeasurableIntegral}(1)(6)可得对任意的$z\in Z$有：
	\begin{equation*}
		\int_{Y}p_{f\mid g}(y\mid z)\dif\mu=1
	\end{equation*}
	对任意的$z\in Z$和$A\in\mathscr{A}$，定义：
	\begin{equation*}
		K_{f\mid g}(z,A)\coloneq\int_{A}p_{f\mid g}(y\mid z)\dif\mu
	\end{equation*}
	对固定的$z$，由\cref{prop:NonnegativeMeasurableIntegral}(2)和\cref{theo:MeasurableCountableIntegral}可知$K_{f\mid g}(z,\cdot)$是$(Y,\mathscr{A})$上的概率测度；对固定的$A$，由\cref{theo:Fubini}(2)可知$K_{f\mid g}(\cdot,A)$关于$\mathscr{B}$可测。因此$K_{f\mid g}$是从$(Z,\mathscr{B})$到$(Y,\mathscr{A})$的概率核。又因为$p_{f\mid g}=q$ a.e.于$(Y\times Z,\mathscr{A}\times\mathscr{B},\mu\times Pg^{-1})$，由\cref{theo:Fubini}(2)和\cref{prop:MeasurableIntegral}(8)可知对任意的$A\in\mathscr{A},\;B\in\mathscr{B}$有：
	\begin{align*}
		&\int_{B}K_{f\mid g}(z,A)\dif Pg^{-1}=\int_{A\times B}p_{f\mid g}(y\mid z)\dif(\mu\times Pg^{-1}) \\
		=&\int_{A\times B}q(y,z)\dif(\mu\times Pg^{-1})=P(f,g)^{-1}(A\times B)=P(\{f\in A\}\cap\{g\in B\})
	\end{align*}
	于是$K_{f\mid g}$是$f$关于$g$的条件概率核，$p_{f\mid g}$是$f$关于$g$相对于$\mu$的条件概率密度函数。
\end{proof}
\begin{property}\label{prop:ConditionalProbabilityDensity}
	设$f,g$分别是从概率空间$(X,\mathscr{F},P)$到可测空间$(Y,\mathscr{A})$和$(Z,\mathscr{B})$的可测映射，$\mu,\nu$分别是$(Y,\mathscr{A})$和$(Z,\mathscr{B})$上的$\sigma$有限测度。若$P(f,g)^{-1}\ll\mu\times\nu$，由\cref{theo:RadonNikodym}和\cref{theo:MarginalDensity}，取处处有限、非负的：
	\begin{equation*}
		p_{f,g}\coloneq\dfrac{\dif P(f,g)^{-1}}{\dif(\mu\times\nu)},\quad p_f\coloneq\dfrac{\dif Pf^{-1}}{\dif\mu},\quad p_g\coloneq\dfrac{\dif Pg^{-1}}{\dif\nu}
	\end{equation*}
	并且它们满足：
	\begin{equation*}
		p_f(y)=\int_{Z}p_{f,g}(y,z)\dif\nu\;\text{a.e.于$(Y,\mathscr{A},\mu)$},\quad
		p_g(z)=\int_{Y}p_{f,g}(y,z)\dif\mu\;\text{a.e.于$(Z,\mathscr{B},\nu)$}
	\end{equation*}
	则：
	\begin{enumerate}
		\item(条件概率密度公式) 定义：
		\begin{equation*}
			p_{f\mid g}(y\mid z)\coloneq
			\begin{cases}
				\dfrac{p_{f,g}(y,z)}{p_g(z)},&p_g(z)>0 \\
				0,&p_g(z)=0
			\end{cases}
		\end{equation*}
		则$p_{f\mid g}$是$f$关于$g$相对于$\mu$的条件概率密度函数；
		\item(全概率公式) 有：
		\begin{equation*}
			p_f(y)=\int_{Z}p_{f\mid g}(y\mid z)p_g(z)\dif\nu\;\text{a.e.于$(Y,\mathscr{A},\mu)$}
		\end{equation*}
		\item(贝叶斯公式) 存在$N\in\mathscr{A}$满足$\mu(N)=0$且对任意的$y\in\{p_f(y)>0\}\setminus N$有：
		\begin{equation*}
			p_{g\mid f}(z\mid y)=\frac{p_{f\mid g}(y\mid z)p_g(z)}{p_f(y)}\;\text{a.e.于$(Z,\mathscr{B},\nu)$}
		\end{equation*}
	\end{enumerate}
\end{property}
\begin{proof}
	(1)记$N=\{z\in Z:p_g(z)=0\}$，由\cref{prop:MeasurableFunction}(1.b)(1.a)和\cref{prop:SigmaField}(4)可知$N\in\mathscr{B}$。根据\cref{theo:MarginalDensity}可得：
	\begin{equation*}
		Pg^{-1}(N)=\int_{N}p_g(z)\dif\nu=0,\quad\int_{Y\times N}p_{f,g}(y,z)\dif(\mu\times\nu)=\int_{N}p_g(z)\dif\nu=0
	\end{equation*}
	可以验证$p_{f\mid g}$是$(Y\times Z,\mathscr{A}\times\mathscr{B})$上的非负可测函数。任取$A\in\mathscr{A},\;B\in\mathscr{B}$，由\cref{lem:IntChangeOfMeasure}、\cref{prop:MeasurableIntegral}(6)、\cref{prop:SigmaField}(4)、\cref{theo:MeasurableCountableIntegral}和\cref{theo:Fubini}(2)可得：
	\begin{align*}
		&\int_{B}\left[\int_{A}p_{f\mid g}(y\mid z)\dif\mu\right]\dif Pg^{-1}=\int_{B}\left[\int_{A}p_{f\mid g}(y\mid z)\dif\mu\right]p_g(z)\dif\nu \\
		=&\int_{B\setminus N}\left[\int_{A}p_{f,g}(y,z)\dif\mu\right]\dif\nu=\int_{A\times B}p_{f,g}(y,z)\dif(\mu\times\nu)=P(f,g)^{-1}(A\times B)
	\end{align*}
	第三式等于第四式用到了\cref{theo:MeasurableCountableIntegral}和\cref{prop:NonnegativeMeasurableIntegral}(2)：
	\begin{equation*}
		0=\int_{Y\times N}p_{f,g}(y,z)\dif(\mu\times\nu)\geqslant\int_{A\times(B\cap N)}p_{f,g}(y,z)\dif(\mu\times\nu)\geqslant0
	\end{equation*}
	对任意的$C\in\mathscr{A}\times\mathscr{B}$，定义：
	\begin{equation*}
		\varphi(C)\coloneq\int_{C}p_{f\mid g}(y\mid z)\dif(\mu\times Pg^{-1})
	\end{equation*}
	由$\varphi(Y\times Z)=1$、\cref{prop:NonnegativeMeasurableIntegral}(2)和\cref{theo:MeasurableCountableIntegral}可知$\varphi$是$(Y\times Z,\mathscr{A}\times\mathscr{B})$上的概率测度，根据\cref{theo:Fubini}可得$\varphi$与$P(f,g)^{-1}$在所有可测矩形上相等。由\cref{lem:PiExtensionMeasure}可得$\varphi=P(f,g)^{-1}$，从而$p_{f\mid g}$是$P(f,g)^{-1}$关于$\mu\times Pg^{-1}$的Radon-Nikodym导数。由\cref{prop:ConditionalProbabilityDensityBasic}(2)可知$p_{f\mid g}$是$f$关于$g$相对于$\mu$的条件概率密度函数。\par
	(2)任取$A\in\mathscr{A}$。由\cref{prop:ConditionalProbabilityDensityBasic}(1)、$Pg^{-1}\ll\nu$、\cref{lem:IntChangeOfMeasure}和\cref{theo:Fubini}(2)可得：
	\begin{align*}
		&Pf^{-1}(A)=P(\{f\in A\}\cap\{g\in Z\})=\int_{Z}\left[\int_{A}p_{f\mid g}(y\mid z)\dif\mu\right]\dif Pg^{-1} \\
		=&\int_{Z}\left[\int_{A}p_{f\mid g}(y\mid z)\dif\mu\right]p_g(z)\dif\nu=\int_{A}\left[\int_{Z}p_{f\mid g}(y\mid z)p_g(z)\dif\nu\right]\dif\mu
	\end{align*}
	由\cref{theo:RadonNikodym}即可得出结论。\par
	(3)由(1)可知对$p_f(y)>0$的$y$有：
	\begin{equation*}
		p_{g\mid f}(z\mid y)=\frac{p_{f,g}(y,z)}{p_f(y)}
	\end{equation*}
	并且：
	\begin{equation*}
		p_{f,g}(y,z)=p_{f\mid g}(y\mid z)p_g(z)\;\text{a.e.于$(Y\times Z,\mathscr{A}\times\mathscr{B},\mu\times\nu)$}
	\end{equation*}
	令：
	\begin{equation*}
		E=\left\{(y,z)\in Y\times Z:p_{f,g}(y,z)\neq p_{f\mid g}(y\mid z)p_g(z)\right\}
	\end{equation*}
	则$(\mu\times\nu)(E)=0$。对任意$y\in Y$，记$E_y=\{z\in Z:(y,z)\in E\}$\info{$E$和$E_y$的可测性与相关函数的可测性}，由\cref{theo:MeasurableCountableIntegral}和\cref{theo:Fubini}(2)可得：
	\begin{align*}
		0&=(\mu\times\nu)(E)=\int_{Y\times Z}I_E(y,z)\dif(\mu\times\nu)=\int_{Y}\left[\int_{Z}I_E(y,z)\dif\nu\right]\dif\mu \\
		&=\int_{Y}\left[\int_{Z}I_{E_y}(z)\dif\nu\right]\dif\mu=\int_{Y}\nu(E_y)\dif\mu
	\end{align*}
	因为$\nu(E_y)\geqslant0$，根据\cref{prop:MeasurableIntegral}(9)可知$\nu(E_y)=0\;$a.e.于$(Y,\mathscr{A},\mu)$。记$N=\{\nu(E_y)\ne0\}$，则$\mu(N)=0$，对于给定的$y\notin N$，$\nu(\{z:p_{f,g}(y,z)\neq p_{f\mid g}(y\mid z)p_g(z)\})=0$，于是对任意的$y\in\{p_f(y)>0\}\setminus N$有：
	\begin{equation*}
		p_{g\mid f}(z\mid y)=\frac{p_{f\mid g}(y\mid z)p_g(z)}{p_f(y)}\;\text{a.e.于$(Z,\mathscr{B},\nu)$}\qedhere
	\end{equation*}
\end{proof}
